\keys{Requirements Engineering; Multi-Agent Systems; Generative Artificial Intelligence; User Stories; Automatic Speech Recognition}

\begin{abstract}{Generative Artificial Intelligence and Software Engineering: Challenges and Opportunities}
% TODO: insert the quantitative numbers (WER/CER, F1, Cronbach's alpha,
% coverage) in the results sentence once the experiments are completed.
Requirements elicitation from meetings is an expensive, manual, and error-prone activity, in which practitioners must simultaneously participate in the discussion and record it as formal artifacts. This work proposes a multi-agent system, based on generative artificial intelligence, that automates this process end to end, transforming meeting audio into structured requirements documentation. The architecture comprises four specialized agents responsible, respectively, for the automatic transcription of the audio, the identification of personas and the generation of user stories following the INVEST criteria, the generation of the requirements document in two selectable formats (the IEEE~830 standard and a client-specific format), and the synthesis of a domain diagram. The agents communicate through a file-based sequential pipeline, which provides loose coupling, fault isolation, and entry-point flexibility. The system was validated in a real-world setting, through a partnership with a public-sector agency, using a manually produced official requirements document as ground truth. The quantitative evaluation covers transcription accuracy, information extraction quality, the document quality as perceived by practitioners, and error propagation across the pipeline. % TODO: Results indicate [insert numbers].
\end{abstract}
